% © 2025 Ing. David Jaroš — CC BY-NC-ND 4.0
%
% This work is licensed under a Creative Commons Attribution-NonCommercial-NoDerivatives
% 4.0 International License (CC BY-NC-ND 4.0).

\documentclass[11pt,a4paper]{article}
\usepackage{amsmath, amssymb, amsthm}
\usepackage{hyperref}
\usepackage{geometry}
\geometry{margin=2.5cm}

\newtheorem{theorem}{Theorem}
\newtheorem{proposition}{Proposition}
\newtheorem{definition}{Definition}
\newtheorem{rule}{Rule}
\newtheorem{remark}{Remark}

\title{Canonical Projection Rules\\
\large Unified Biquaternion Theory — Canonical Theory Series}
\author{Ing.\ David Jaroš}
\date{2025}

\begin{document}
\maketitle

\begin{abstract}
The biquaternionic field $\Theta$ encodes geometry, gauge fields, and matter
simultaneously. Physical observables are extracted from $\Theta$ through
structured projection rules. This document defines all canonical projection
rules used in UBT. It does not contain new derivations; it collects and
cross-references existing definitions from canonical source files.
\end{abstract}

\section{Overview}

A biquaternion $b \in \mathcal{B} = \mathbb{H}\otimes_\mathbb{R}\mathbb{C}$
has the decomposition:
\begin{equation}
  b \;=\; b_0 + b_1\mathbf{i} + b_2\mathbf{j} + b_3\mathbf{k},
  \qquad b_\mu \in \mathbb{C},
\end{equation}
where $\{\mathbf{i},\mathbf{j},\mathbf{k}\}$ are the quaternion units.
Each $b_\mu = \mathrm{Re}(b_\mu) + i\,\mathrm{Im}(b_\mu)$, giving 8 real
degrees of freedom per spacetime point for a single biquaternion.

\noindent\textbf{Source}: \texttt{canonical/fields/biquaternion\_algebra.tex}, \S1.

\section{Rule P1: Metric Projection}

\begin{rule}[Metric Extraction]\label{rule:metric}
The classical spacetime metric is extracted by:
\begin{equation}\label{eq:P1}
  g_{\mu\nu}
  \;=\;
  \mathrm{Re}\bigl(\mathcal{G}_{\mu\nu}[\Theta]\bigr)
  \;=\;
  \mathrm{Re}\,\mathrm{Tr}\bigl(\partial_\mu\Theta\cdot(\partial_\nu\Theta)^\dagger\bigr).
\end{equation}
\textbf{Forbidden}: introducing $g_{\mu\nu}$ as an independent variable without
explicit reference to \eqref{eq:P1}.
\end{rule}

\noindent\textbf{Status}: PROVED [L0] (algebraic identity).\\
\noindent\textbf{Source}: \texttt{canonical/geometry/metric.tex}, Rule M1;
\texttt{canonical/geometry/biquaternion\_metric.tex}, \S2.

\section{Rule P2: Phase Projection}

\begin{rule}[Phase Frame Projection]\label{rule:phase}
For a phase angle $\varphi \in [0,2\pi)$, the phase-projected metric is:
\begin{equation}
  g_{\mu\nu}^{(\varphi)}
  \;=\;
  \mathrm{Re}\bigl(e^{-i\varphi}\,\mathcal{G}_{\mu\nu}[\Theta]\bigr).
\end{equation}
The standard choice $\varphi = 0$ (Rule P1) is a gauge choice.
Different phase frames are connected by a U(1) phase transformation.
\end{rule}

\noindent\textbf{Status}: PROVED [L0].\\
\noindent\textbf{Source}: \texttt{canonical/geometry/phase\_projection.tex}, \S2.

\section{Rule P3: Gauge Field Extraction}

\begin{rule}[Electromagnetic Potential]\label{rule:A}
The U(1) electromagnetic potential is extracted from the phase of $\Theta$:
\begin{equation}
  A_\mu(x) \;=\; \frac{1}{e}\,\partial_\mu\,\mathrm{arg}(\Theta_{\text{scalar}}),
\end{equation}
where $\mathrm{arg}(\Theta_{\text{scalar}})$ is the complex argument of the
scalar component of $\Theta$.
More precisely, $A_\mu$ arises from the U(1) connection in the covariant
derivative $D_\mu = \partial_\mu - ieA_\mu$.
\end{rule}

\begin{rule}[Non-Abelian Gauge Fields]\label{rule:YM}
The SU(3) color gauge field $G_\mu^a$ and SU(2)$_L$ weak gauge field $W_\mu^a$
are extracted from the quaternionic components of $\Theta$ under the respective
group actions:
\begin{itemize}
  \item $G_\mu^a$: from left SU(3) action on the $3\times 3$ color block of $\Theta$
  \item $W_\mu^a$: from left SU(2) action on the spinorial component of $\Theta$
  \item $B_\mu$: from the right U(1)$_Y$ scalar phase (overlaps with Rule P3)
\end{itemize}
\end{rule}

\noindent\textbf{Status}: PROVED [L0] for group structure; full extraction
map is partial [L2].\\
\noindent\textbf{Source}: \texttt{canonical/interactions/sm\_gauge.tex}, \S2--\S4;
\texttt{canonical/bridges/gauge\_emergence\_bridge.tex}.

\section{Rule P4: Matter Field Extraction}

\begin{rule}[Fermion Fields]\label{rule:fermion}
Spinorial matter fields are identified with the off-diagonal biquaternionic
components of $\Theta$ in the 4-spinor representation:
\begin{equation}
  \psi_\alpha(x) \;\subset\; \Theta(x,\tau)\big|_{\text{spinor block}}.
\end{equation}
The Dirac spinor $\psi$ satisfies $(i\gamma^\mu D_\mu - m)\psi = 0$
in the appropriate limit.
\end{rule}

\noindent\textbf{Status}: PRESENT [P] — identification is motivated;
complete projection map is partial.\\
\noindent\textbf{Source}: \texttt{canonical/interactions/qed.tex}, \S4.

\section{Rule P5: Stress-Energy Extraction}

\begin{rule}[Physical Stress-Energy]\label{rule:Tmunu}
The physical stress-energy tensor is extracted by Hilbert variation:
\begin{equation}
  T_{\mu\nu}
  \;=\;
  -\frac{2}{\sqrt{-g}}\,
  \frac{\delta\bigl(\sqrt{-g}\,\mathcal{L}_{\mathrm{matter}}\bigr)}{\delta g^{\mu\nu}},
\end{equation}
where $\mathcal{L}_{\mathrm{matter}} = \mathrm{Re}\,\mathrm{Tr}[(D_\mu\Theta)^\dagger(D^\mu\Theta)]$
is the matter Lagrangian density extracted from the real part of the
biquaternionic kinetic term.
\end{rule}

\noindent\textbf{Status}: PROVED [L1].\\
\noindent\textbf{Source}: \texttt{canonical/geometry/stress\_energy.tex};
\texttt{consolidation\_project/GR\_closure/step4\_offshell\_Tmunu.tex}.

\section{Rule P6: Observable Limit}

\begin{rule}[Classical Observable Limit]\label{rule:classical}
All classical observables are recovered in the joint limit:
\begin{enumerate}
  \item Real-sector projection: $\psi \to 0$ (or equivalently $\varphi = 0$),
  \item Weak-field limit for GR (or strong field — both are covered),
  \item Classical matter limit: $\hbar \to 0$ for particle trajectories.
\end{enumerate}
The biquaternionic corrections (imaginary components) are invisible to
classical measurements because ordinary matter couples only to $g_{\mu\nu}$,
not to the full $\mathcal{G}_{\mu\nu}$.
\end{rule}

\noindent\textbf{Status}: PROVED [L1] for gravitational sector; PRESENT [P]
for quantum sectors.\\
\noindent\textbf{Source}: \texttt{canonical/geometry/gr\_as\_limit.tex}, \S2;
\texttt{THEORY/README.md}, \S3.

\section{Summary of Projection Rules}

\begin{center}
\begin{tabular}{lllll}
\hline
\textbf{Rule} & \textbf{Input} & \textbf{Output} & \textbf{Status} & \textbf{Source} \\
\hline
P1 & $\Theta$ & $g_{\mu\nu}$ (metric) & [L0] & \texttt{metric.tex} M1 \\
P2 & $\mathcal{G}_{\mu\nu}$, $\varphi$ & $g_{\mu\nu}^{(\varphi)}$ & [L0] & \texttt{phase\_projection.tex} \\
P3 & phase of $\Theta$ & $A_\mu$, $B_\mu$ & [L0]/[L2] & \texttt{sm\_gauge.tex} \\
P3b & quaternion of $\Theta$ & $G_\mu^a$, $W_\mu^a$ & [L0]/[L2] & \texttt{sm\_gauge.tex} \\
P4 & spinor block & $\psi_\alpha$ & [P] & \texttt{qed.tex} \S4 \\
P5 & $\mathcal{L}_{\mathrm{matter}}$ & $T_{\mu\nu}$ & [L1] & \texttt{stress\_energy.tex} \\
P6 & $\Theta$ + limits & classical fields & [L1] & \texttt{gr\_as\_limit.tex} \\
\hline
\end{tabular}
\end{center}

\section{Cross-References}

\begin{itemize}
  \item Axioms: \texttt{THEORY/canonical/canonical\_axioms.tex}
  \item Action: \texttt{THEORY/canonical/canonical\_action.tex}
  \item Metric: \texttt{THEORY/canonical/canonical\_metric\_definition.tex}
  \item Field equations: \texttt{THEORY/canonical/canonical\_field\_equations.tex}
  \item Full algebra: \texttt{canonical/fields/biquaternion\_algebra.tex}
  \item Phase projection: \texttt{canonical/geometry/phase\_projection.tex}
  \item Gauge groups: \texttt{canonical/bridges/gauge\_emergence\_bridge.tex}
\end{itemize}

\end{document}
